\documentclass[a4paper,12pt]{article}
\usepackage[utf8]{inputenc}
\usepackage[italian]{babel}
\usepackage{listings}
\usepackage{xcolor}
\usepackage{float}
\usepackage{placeins}
\usepackage{url}
\setlength{\emergencystretch}{3em}

\lstset{
  language=C,
  backgroundcolor=\color{lightgray!20},
  frame=single,
  breaklines=true,
  numbers=left,
  numberstyle=\tiny,
  stepnumber=1,
  numbersep=5pt,
  captionpos=b,
  basicstyle=\ttfamily\small,
  keywordstyle=\color{blue}\bfseries,
  commentstyle=\color{green!60!black},
  stringstyle=\color{red},
  showstringspaces=false
}

\title{LabReSiD26 -- Hands-On 13}
\author{Davide Ferrara -- 518629}
\date{}

\begin{document}

\maketitle

% ---------------------------------------------------------------------------
\section{Esercizio 1: Somma parallela di un array}
% ---------------------------------------------------------------------------

L'esercizio consiste nel dividere un array di $N = 1\,000\,000$ interi in
$T = 4$ segmenti uguali, calcolare la somma parziale di ciascun segmento
in un thread separato tramite \texttt{pthread\_create()}, raccogliere i
risultati con \texttt{pthread\_join()} e verificare che la somma totale
coincida con quella calcolata sequenzialmente.

\subsection*{Struttura degli argomenti}

Ogni thread riceve un puntatore a una struttura \texttt{worker\_args} che
descrive il segmento di lavoro e dove scrivere il risultato parziale.
Usare una struttura è necessario perché \texttt{pthread\_create()} accetta
un unico argomento \texttt{void *}.

\begin{lstlisting}[caption={Struttura argomento thread e funzione job}]
/* Struttura dati del worker*/
typedef struct {
  const int *array;
  int from;
  int to;
  ssize_t result;
} worker_args;

/* Job del thread, fa il casting dell'args da void* in worker_args*
 * in questo modo il programma sa come leggere la struttura in memoria.
 * Uso la funziona gettid() per ottenere il thread id.
 * Il thread chiama partial_sum() e salva il risultato nella struttura,
 * il thread ritorna NULL.
 * */
static void job(void *args) {
  worker_args *w = (worker_args *)args;
  pid_t tid = gettid();

  w->result = partial_sum(w->array, w->from, w->to);
  printf("[SUM.c][Thread %d] Somma parziale: %lu\n", tid, w->result);

  pthread_exit(NULL);
}
\end{lstlisting}

\subsection*{Creazione e raccolta dei thread}

I loop di creazione e raccolta sono separati: il primo lancia tutti i
thread in parallelo, il secondo raccoglie i risultati con
\texttt{pthread\_join()}.
I segmenti dell'array sono disgiunti, quindi i thread non condividono
mai dati in scrittura: non serve nessun meccanismo di sincronizzazione.
\texttt{pthread\_join()} attende la terminazione del thread e raccoglie
il valore restituito da \texttt{return} o \texttt{pthread\_exit()}.

\begin{lstlisting}[caption={Loop di creazione e raccolta separati}]
worker_args w[T] = {0};
pthread_t tid[T];

/* Stabilisco lo step per dividere l'array in base a T */
int step = SIZE / T;
int from = 0, to = step;

/* Lancio tutti i thread in parallelo: segmenti disgiunti,
 * nessuna sincronizzazione necessaria. */
for (int i = 0; i < T; i++) {
  w[i].array = a;
  w[i].from = from;
  w[i].to = to;
  from = to;
  to += step;
  pthread_create(&tid[i], NULL, (void *)&job, (void *)&w[i]);
}

long long psum = 0;
for (int i = 0; i < T; i++) {
  if (pthread_join(tid[i], NULL) == 0)
    psum += w[i].result;
}
\end{lstlisting}

\begin{lstlisting}[language=bash, caption={Test Esercizio 1}]
$ ./sum
[SUM.c] Somma sequenziale: 49500000
[SUM.c][Thread 29728] Somma parziale: 12375000
[SUM.c][Thread 29729] Somma parziale: 12375000
[SUM.c][Thread 29730] Somma parziale: 12375000
[SUM.c][Thread 29731] Somma parziale: 12375000
[SUM.c] Somma parallela (T=4): 49500000
[OK] i risultati coincidono!
\end{lstlisting}

% ---------------------------------------------------------------------------
\section{Esercizio 2: Barrier -- map-reduce parallelo}
% ---------------------------------------------------------------------------

L'esercizio implementa uno schema map-reduce in due fasi usando
\texttt{pthread\_barrier\_t}: $T$ thread eseguono la fase di \emph{map}
in parallelo (calcolo del quadrato di ogni elemento nel proprio segmento),
si sincronizzano alla barrier, poi il solo thread che riceve
\texttt{PTHREAD\_BARRIER\_SERIAL\_THREAD} esegue la fase di \emph{reduce}
(somma di tutti i quadrati).

\subsection*{Inizializzazione della barrier}

La barrier va creata nel \texttt{main} prima di lanciare i thread,
specificando il numero di thread che devono raggiungerla prima che
vengano sbloccati:

\begin{lstlisting}[caption={Inizializzazione e distruzione della barrier}]
pthread_barrier_init(&barrier, NULL, T);
/* ... creazione e join dei thread ... */
pthread_barrier_destroy(&barrier);
\end{lstlisting}

\subsection*{Fase map e sincronizzazione}

Ogni thread eleva al quadrato gli elementi nel proprio segmento, poi
chiama \texttt{pthread\_barrier\_wait()}.
La barrier blocca ogni thread finché tutti e $T$ non l'hanno raggiunta,
garantendo che il reduce non inizi prima che il map sia completato da tutti.
Lo standard POSIX garantisce che esattamente uno tra i $T$ thread riceva
\texttt{PTHREAD\_BARRIER\_SERIAL\_THREAD}: è l'unico meccanismo di
selezione ``uno tra N'' offerto da pthreads senza mutex.

\begin{lstlisting}[caption={Funzione job: map, barrier e reduce}]
static void job(void *args) {
  worker_args *w = (worker_args *)args;
  pid_t tid = gettid();

  printf("[Thread %d] Iniziato il map...\n", tid);
  map(w->array, w->from, w->to);

  int rc = pthread_barrier_wait(&barrier);
  if (rc == PTHREAD_BARRIER_SERIAL_THREAD) {
    printf("[Thread %d Seriale %d] Iniziato il reduce...\n", tid, rc);
    w->result = reduce(w->array);
  }

  pthread_exit(NULL);
}
\end{lstlisting}

\medskip
\noindent\textbf{Nota:} le slide riportano
332\,833\,500\,000, che corrisponde a $N = 1\,000\,000$ elementi.
Con $N = 800\,000$ (come specificato nell'esercizio) il valore corretto è
peró $266\,266\,800\,000$.

\begin{lstlisting}[language=bash, caption={Test Esercizio 2}]
$ ./mapreduce
Somma dei quadrati (sequenziale): 266266800000
[Thread 58309] Iniziato il map...
[Thread 58310] Iniziato il map...
[Thread 58311] Iniziato il map...
[Thread 58312] Iniziato il map...
[Thread 58311 Seriale -1] Iniziato il reduce...
Somma dei quadrati (parallela): 266266800000
[OK] i risultati coincidono!
\end{lstlisting}

\end{document}
